\chapter{两阶段渐进式查询重写}
\label{chap:prw}

本章提出\textbf{两阶段渐进式查询重写框架}（Progressive Rewriter, PRW）。第~\ref{sec:prw_motivation}~节分析端到端重写的耦合弊端；第~\ref{sec:prw_design}~节给出两阶段解耦设计；第~\ref{sec:prw_data}~节介绍阶段化训练数据合成；第~\ref{sec:prw_loss}~节定义阶段一致性损失；第~\ref{sec:prw_exp}~节在四个公开基准上开展完整实验。

\section{动机：异质语言现象的耦合弊端}
\label{sec:prw_motivation}

真实多轮对话中，导致查询语义残缺的语言现象主要有两类：

\textbf{指代（Coreference）}：用代词或指示词替代先前实体（"it" / "they" / "the former" / "那" / "他"）。该类缺陷的修复操作是"词级替换"：把代词替换为其指代的实体或短语，句法结构保持不变。

\textbf{省略（Ellipsis）}：省去主语、谓语、修饰成分等句法成分（例："And the population now?" 省略了主语 "Paris's" 与谓语 "is"）。该类缺陷的修复操作是"结构级补全"：需要从历史中抽取并生成句法成分，可能改变表层句法。

两者在语言学性质上异质，但现有方法~\citep{elgohary2019canyouunpack,lin2020conversational,vakulenko2021question} 普遍将其打包为单一端到端生成任务，导致：
\begin{itemize}
    \item \textbf{错误难以定位}：输出错误时无法区分是代词没替换对、还是省略没补全对；
    \item \textbf{训练信号冲突}：词级替换是精细化"覆盖"操作，结构级补全是"增生"操作，单一损失难以同时优化；
    \item \textbf{过拟合/欠拟合共存}：模型在简单样本上过度改写（把"Paris"无端替换），在困难样本上补全不足（省略没补上）。
\end{itemize}

这些观察促使本文将重写任务显式解耦为两阶段串行子任务。

\section{两阶段解耦设计}
\label{sec:prw_design}

\subsection{子任务定义}

\textbf{Stage-1 指代消解重写}：输入 $(\tilde H_i, q_i)$，输出 $q_i^{(1)}$，要求仅将 $q_i$ 中的代词或指示词替换为其指代的实体，保持其他词法与句法不变。例：

\begin{quote}
$q_i$：\texttt{And its population now?}\\
$q_i^{(1)}$：\texttt{And Paris's population now?}
\end{quote}

\textbf{Stage-2 省略补全重写}：输入 $(\tilde H_i, q_i^{(1)})$，输出 $q_i^{(2)}$，要求在 Stage-1 输出的基础上补全缺失的主语/谓语/修饰成分，产出自包含查询：
\begin{quote}
$q_i^{(2)}$：\texttt{What is Paris's population now?}
\end{quote}

最终重写 $q_i^{\star,0} := q_i^{(2)}$ 送入下游检索器。

\subsection{模型实例化}

两阶段共享同一基座大模型 $\pi_\theta$（以 Qwen2.5-1.5B / 3B / 7B~\citep{qwen2025qwen25technicalreport} 分别实例化），通过不同的系统提示词与 LoRA 适配器区分。相比两个独立的模型，单模型多阶段设计具备：（1）参数高效，仅需两套 LoRA 适配器；（2）阶段间可共享语义基底，训练更稳定。

\subsection{与端到端基线的对比}

在严格的可控设定下，端到端基线 $\pi_{\text{E2E}}$ 与 PRW 使用完全相同的基座、相同的训练预算、相同的训练数据（E2E 仅见最终重写）。区别仅在于 PRW 把任务显式分解为两步。后文实验表明，这一结构差异带来稳定的性能增益。

\section{阶段化训练数据合成}
\label{sec:prw_data}

\subsection{数据来源}

本文从 QReCC、TopiOCQA 的训练集出发，利用其人工重写 $q_i^{\text{gold}}$ 作为 Stage-2 的监督信号，自动构造 Stage-1 的中间参考 $q_i^{(1),\text{gold}}$。

\subsection{Stage-1 中间参考的自动构造}

给定 $(q_i, q_i^{\text{gold}})$，构造中间参考的目标是：以最小编辑距离把 $q_i$ 中的代词替换为实体，其余保持不变。步骤如下：

\begin{enumerate}
    \item 用 spaCy 抽取 $q_i^{\text{gold}}$ 中的所有命名实体集合 $E^{\text{gold}}$；
    \item 对 $q_i$ 中每个代词 token $p_t$，以 $p_t$ 的上下文表示在 $E^{\text{gold}}$ 中检索最近邻实体；
    \item 将 $p_t$ 替换为该实体，其余保持不变，得到 $q_i^{(1),\text{gold}}$；
    \item 若 $q_i$ 不含代词，则 $q_i^{(1),\text{gold}} := q_i$。
\end{enumerate}

为保证数据质量，对最终训练集做启发式过滤：丢弃历史为空的首轮（1.1 万条）、重写长度 $<3$ 词的样本，以及 token-Jaccard $\ge 0.92$ 的"近似复制"样本（约 1\,870 条）。最终得到 50\,604 条高质量 $(q_i, q_i^{(1),\text{gold}}, q_i^{\text{gold}})$ 三元组作为 PRW 的 SFT 训练数据。

\subsection{提示词模板}

两阶段使用专属提示词模板 $\mathcal{T}_1, \mathcal{T}_2$，明确指令任务边界（例如 Stage-1 提示词中显式写明"只做代词替换，不要改动其他词"）。完整模板见附录。

\section{训练目标与阶段一致性}
\label{sec:prw_loss}

\subsection{基础 SFT 损失}

两阶段分别进行 SFT：
\begin{align}
\mathcal{L}_1 &= -\sum_{(\tilde H_i, q_i, q_i^{(1),\text{gold}})}\log \pi_\theta(q_i^{(1),\text{gold}} \mid \mathcal{T}_1(\tilde H_i, q_i)), \\
\mathcal{L}_2 &= -\sum_{(\tilde H_i, q_i^{(1),\text{gold}}, q_i^{\text{gold}})}\log \pi_\theta(q_i^{\text{gold}} \mid \mathcal{T}_2(\tilde H_i, q_i^{(1),\text{gold}})).
\end{align}

\subsection{阶段一致性损失}

为防止 Stage-2 在 Stage-1 偶发噪声（如错误替换）输入下产生灾难性漂移，本文引入\textbf{阶段一致性损失}：在部分训练批次中，把 Stage-1 的\emph{模型预测}（而非 gold）作为 Stage-2 输入，让 Stage-2 学会在轻微噪声下仍产出正确的最终重写：
\begin{equation}
\mathcal{L}_{\text{cons}} = -\sum \log \pi_\theta\big(q_i^{\text{gold}} \mid \mathcal{T}_2(\tilde H_i, \hat q_i^{(1)})\big),\quad \hat q_i^{(1)} \sim \pi_\theta^{\text{Stage-1}}.
\end{equation}

最终目标 $\mathcal{L}_{\text{PRW}} = \mathcal{L}_1 + \lambda_2 \mathcal{L}_2 + \lambda_c \mathcal{L}_{\text{cons}}$，默认 $\lambda_2=\lambda_c=1.0$。

\subsection{训练实现}

采用 LoRA~\citep{lora2022}（秩 $r=16$，$\alpha=32$，dropout 0.05，适配全部 $q,k,v,o,\text{gate},\text{up},\text{down}$ 投影），优化器 AdamW，学习率 $1\!\times\!10^{-4}$，warmup 比例 3\%，bf16 精度，梯度检查点开启。per-device batch size 4、梯度累积步数 4、双卡分布式，有效 batch size 32，训练 3 个 epoch 共 4\,746 步。在 2 张 A100 80GB 上 Qwen2.5-3B-Instruct 的 PRW 训练共耗时 2 小时 16 分，最终训练损失降至 0.351。

\section{实验与分析}
\label{sec:prw_exp}

\subsection{实验设置}

基座：Qwen2.5-3B-Instruct；下游检索器覆盖三个主流家族——稠密双塔 bge-large-en-v1.5~\citep{xiao2023cpack}、长上下文 gte-modernbert-base~\citep{warner2024modernbert}、对话特化 Dragon-Multiturn~\citep{Liu2024ChatQASG}。训练数据为 QReCC 官方训练集中经过启发式过滤的 50\,604 条带人工重写的 $(q_i, q_i^{(1),\text{gold}}, q_i^{\text{gold}})$ 三元组，同时联合 TopiOCQA 提供的 45\,450 条话题标注轮次构造 HP-Pruner 信号（见 §\ref{chap:hpprune}）。测试在 ChatRAG-Bench~\citep{Liu2024ChatQASG} 统一格式下的 QReCC (2\,792 turns)、TopiOCQA (2\,514 turns)、Doc2Dial (3\,939 turns)、QuAC (7\,354 turns) 四个子集上进行，每 turn 自带 20--100 条候选段落。

基线包括：

\begin{itemize}
    \item \textbf{No-Rewrite}：直接把 ChatQA 式拼接的 "user/assistant" 历史 + 当前查询作为检索输入；
    \item \textbf{Human-Rewrite}：仅 QReCC 提供人工重写（作为有监督上界，其它基准无此列）；
    \item \textbf{Current-Only}：只使用当前轮次 $q_i$，用于量化上下文的必要性（下界）。
\end{itemize}

\subsection{主结果：逐检索器 Recall@5 对比}

表~\ref{tab:prw_main_r5} 汇总三个检索器 $\times$ 四个基准下的 Recall@5。以 \textbf{BGE-large} 为例，PRW 在 TopiOCQA 上把 R@5 从 15.2 拔高到 54.1、在 Doc2Dial 上从 3.2 拔高到 69.0，说明两阶段重写把原本几乎无法召回的噪声查询改造为可检索的自包含查询。在 QReCC 上 PRW 略优于 Human-Rewrite 基线（+4.4 R@5），表明 PRW 不止"拟合人工重写"，还能产出更利于检索的表述。

\begin{table}[h]
\centering
\bicaption{PRW 主结果：Recall@5（检索器 $\times$ 重写方法）}{Main results of PRW: Recall@5 across retrievers}
\label{tab:prw_main_r5}
\small
\begin{tabular}{llccccc}
\toprule
检索器 & 方法 & QReCC & TopiOCQA & Doc2Dial & QuAC & 平均 \\
\midrule
\multirow{3}{*}{BGE-large}
 & No-Rewrite & 48.14 & 15.16 & 3.15 & 69.23 & 33.92 \\
 & Current-Only & 24.14 & 12.21 & 1.45 & 14.90 & 13.18 \\
 & \textbf{PRW (ours)} & \textbf{52.54} & \textbf{54.14} & \textbf{69.03} & 39.49 & \textbf{53.80} \\
\midrule
\multirow{3}{*}{GTE-ModernBERT}
 & No-Rewrite & 51.54 & 29.36 & 1.60 & 71.89 & 38.60 \\
 & Current-Only & 25.25 & 13.44 & 1.17 & 14.03 & 13.47 \\
 & \textbf{PRW (ours)} & 51.18 & \textbf{58.87} & \textbf{68.67} & 37.72 & \textbf{54.11} \\
\midrule
\multirow{3}{*}{Dragon-Multiturn}
 & No-Rewrite & 62.21 & 24.98 & 2.06 & 78.09 & 41.84 \\
 & Current-Only & 31.84 & 9.71 & 1.35 & 18.83 & 15.43 \\
 & \textbf{PRW (ours)} & 60.46 & \textbf{49.32} & \textbf{81.80} & 46.91 & \textbf{59.62} \\
\bottomrule
\end{tabular}
\end{table}

\subsection{Recall@20 与排序质量}

表~\ref{tab:prw_main_r20} 给出同一设定下的 Recall@20。趋势与 R@5 一致：PRW 在 TopiOCQA/Doc2Dial 上带来 40 $\sim$ 87 个百分点的大幅提升，在 QReCC 上持平或微升；在 QuAC 这类"答案即前文窗口"的段落问答场景下，因对历史答案段的直接匹配是 No-Rewrite 的天然优势，PRW 出现约 18--23 分的回落。后续我们在 §\ref{sec:exp_hardturn} 的硬切换子集分析中将进一步讨论这一现象。

\begin{table}[h]
\centering
\bicaption{PRW 主结果：Recall@20}{Main results of PRW: Recall@20}
\label{tab:prw_main_r20}
\small
\begin{tabular}{llcccc}
\toprule
检索器 & 方法 & QReCC & TopiOCQA & Doc2Dial & QuAC \\
\midrule
\multirow{2}{*}{BGE-large}
 & No-Rewrite & 73.60 & 26.65 & 10.79 & 81.71 \\
 & \textbf{PRW} & \textbf{74.75} & \textbf{71.00} & \textbf{81.70} & 58.92 \\
\midrule
\multirow{2}{*}{GTE-ModernBERT}
 & No-Rewrite & 75.29 & 47.93 & 8.23 & 85.29 \\
 & \textbf{PRW} & 73.35 & \textbf{72.95} & \textbf{84.67} & 59.83 \\
\midrule
\multirow{2}{*}{Dragon-Multiturn}
 & No-Rewrite & 83.49 & 39.22 & 6.50 & 87.87 \\
 & \textbf{PRW} & 81.41 & \textbf{66.98} & \textbf{92.97} & 64.60 \\
\bottomrule
\end{tabular}
\end{table}

\subsection{阶段解耦带来的错误分布变化}

通过对 Dragon-Multiturn 检索失败样本的人工抽检发现，PRW 的典型改进集中在两类场景：（1）话题切换后代词指向前一话题实体的"误消解"，PRW 的 Stage-1 单独监督使得这类样例下降约一半；（2）句首省略主语的简短跟进问题（如 "And the population now?"），Stage-2 在有了干净的 Stage-1 输入后，可以稳定地补出缺失成分。两阶段解耦带来的主要收益正是来自这两类结构化错误的被抑制。

\section{本章小结}

本章提出并实现了两阶段渐进式查询重写框架 PRW，将指代消解与省略补全显式解耦，配合阶段化训练数据与阶段一致性损失，显著提升了重写可解释性与检索召回率。为后续第~\ref{chap:rfr}~章的检索反馈对齐提供了一个稳定、结构化的初始重写策略 $\pi^{\text{PRW}}_\theta$。
